\ifthenelse{\equal{\doclanguage}{English}}
{\chapter{Introduction}}
{\chapter{Introduzione}}
\label{sec:introduction}

\section{Contestualizzazione}
Negli ultimi decenni la quantità di dati prodotti e consumati a livello globale è cresciuta a ritmi senza precedenti. Secondo le stime di Statista, nel 2024 il volume globale delle informazioni create e consumate supererà i 149 zettabyte ($10^{21}$ byte) \cite{statista2024data}, confermando una crescita esponenziale del volume dei dati prodotti a livello mondiale.

Diversi fattori strutturali spiegano questa esplosione: la diffusione dell'Internet of Things (IoT), la trasformazione digitale della sanità e l'emergere delle tecnologie di \textit{Next-Generation Sequencing} (NGS). Il dato non rappresenta più soltanto un sottoprodotto dell'attività digitale, ma una risorsa strategica che richiede nuovi strumenti di gestione e analisi. Si parla del fenomeno dei \textbf{Big Data}, descritto attraverso il concetto delle ``cinque V'': il \textit{volume} massiccio dei dataset, la \textit{velocità} con cui i flussi vengono generati e analizzati, la \textit{varietà} delle informazioni, la \textit{veridicità} della fonte e il \textit{valore} dell'analisi \cite{atzeni2018basi}. In bioinformatica, questo impatto è ancora più evidente. La ricerca farmacologica e lo studio delle reti metaboliche richiedono l'integrazione di dati molto vari, dove il legame tra le entità biologiche assume un'importanza informatica almeno pari a quella delle entità stesse.

\section{Definizione del problema}
Per decenni, la gestione dei dati è stata dominata dai sistemi \textbf{RDBMS} (\textit{Relational Database Management System}), basati su un modello rigorosamente strutturato e su una teoria consolidata che garantisce consistenza e integrità delle informazioni. Tuttavia, mostra limiti evidenti quando viene applicato a domini caratterizzati da forte interconnessione, come le reti biologiche e biomediche.

In tali scenari, le entità non possono essere considerate in modo isolato, ma devono essere interpretate all'interno di una fitta rete di relazioni. Quando una query richiede l'attraversamento di molteplici livelli di connessione, il modello relazionale è costretto a ricorrere a un numero elevato di operazioni di \textit{JOIN}. Questo comportamento introduce un overhead computazionale significativo, che cresce rapidamente con la profondità della query e con la densità della rete, compromettendo la scalabilità delle interrogazioni.

Per far fronte a queste limitazioni, si è progressivamente affermato il paradigma \textbf{NoSQL} (\textit{``Not Only SQL''}) \cite{fowler2015nosql}, termine coniato da Carlo Strozzi nel 1998 per indicare una famiglia eterogenea di sistemi di gestione dei dati caratterizzati da flessibilità dello schema (\textit{schema-agnostic}), scalabilità orizzontale e capacità di distribuzione su infrastrutture commodity. Tale paradigma non sostituisce il modello relazionale, ma lo affianca in contesti in cui la rigidità strutturale rappresenta un vincolo.

All'interno dell'ecosistema NoSQL, i database a grafo hanno assunto un ruolo centrale nella modellazione di sistemi complessi. Essi consentono di rappresentare in modo naturale entità e relazioni, preservando la struttura intrinseca del dominio: i nodi modellano le entità biologiche, mentre gli archi rappresentano le interazioni tra esse. Questa rappresentazione si adatta particolarmente bene alla bioinformatica, dove il valore informativo risiede spesso nelle connessioni piuttosto che nei singoli elementi.

Neo4j rappresenta oggi uno dei database a grafo nativi più diffusi e maturi. Attraverso la libreria \textit{Graph Data Science} (GDS) \cite{neo4j_gds}, il sistema mette a disposizione algoritmi scalabili per l'analisi dei grafi e per l'apprendimento automatico su strutture complesse, integrando in un unico ecosistema storage, query engine e pipeline di analisi.

\section{Obiettivi del lavoro svolto}
Il presente lavoro propone un'analisi congiunta teorica e sperimentale basata sul knowledge graph biomedico PharMeBiNet V3 \cite{konigs2022pharmebinet}, una rete eterogenea di grandi dimensioni composta da circa 8 milioni di nodi e 53 milioni di relazioni, provenienti da molteplici database farmacologici e biologici. La complessità strutturale del grafo lo rende un banco di prova particolarmente rappresentativo per lo studio di algoritmi di Graph Machine Learning in contesti reali.

Il task considerato è la \textit{link prediction} tra nodi di tipo \textit{Chemical} e \textit{Disease}, problematica centrale nel dominio del \textit{drug repurposing}. L'obiettivo consiste nell'identificare potenziali relazioni terapeutiche non ancora documentate, sfruttando la struttura del grafo e le informazioni latenti codificate nelle sue connessioni.

L'analisi si concentra sul confronto tra tre algoritmi: GraphSAGE \cite{hamilton2017inductive}, TransE \cite{bordes2013translating} e HashGNN \cite{neo4j_hashgnn}. Tali modelli sono stati selezionati in quanto rappresentano tre paradigmi distinti di apprendimento su grafi: l'aggregazione di informazioni nei vicini (GraphSAGE), la modellazione geometrica delle relazioni nello spazio latente (TransE) e la gestione esplicita dell'eterogeneità tramite rappresentazioni differenziate per tipo di nodo (HashGNN). Questa diversità metodologica consente di valutare in modo più completo l'impatto delle scelte architetturali sul problema considerato.

Dal punto di vista infrastrutturale, il lavoro prevede la progettazione e l'implementazione di un cluster Neo4j Enterprise distribuito su quattro macchine virtuali in ambiente Azure, configurato secondo un'architettura primary-secondary. L'obiettivo è analizzare l'impatto della scalabilità orizzontale sui tempi di estrazione dei dati necessari alla pipeline di training basata su PyTorch Geometric, variando il numero di nodi attivi nel cluster.

In questo contesto, la valutazione non si limita alla sola accuratezza predittiva dei modelli, ma estende l'analisi alle prestazioni end-to-end della pipeline, con l'obiettivo di quantificare il reale contributo dell'infrastruttura distribuita in scenari di data engineering su grafi biomedici di grande scala.